
\documentclass[10pt,a4paper]{article}

\usepackage[
  margin=1.8cm,
  top=1.55cm,
  bottom=1.55cm
]{geometry}


\usepackage{
  fontspec,
  xeCJK,
  xcolor,
  enumitem,
  titlesec,
  hyperref,
  tabularx,
  parskip
}

% English and Latin characters
\setmainfont[
  Ligatures=TeX
]{TeX Gyre Termes}

% Japanese serif font
\setCJKmainfont[
  BoldFont=HaranoAjiMincho-Bold.otf,
  AutoFakeSlant=true
]{HaranoAjiMincho-Regular.otf}

% Japanese sans-serif font
\setCJKsansfont[
  BoldFont=HaranoAjiGothic-Bold.otf
]{HaranoAjiGothic-Regular.otf}


\definecolor{navy}{HTML}{183650}
\hypersetup{
  colorlinks=true,
  urlcolor=navy
}

\setlength{\parindent}{0pt}
\setlength{\parskip}{3pt}

\setlist[itemize]{
  leftmargin=1.35em,
  topsep=2pt,
  itemsep=1pt,
  parsep=0pt
}

\titleformat{\section}
  {\large\bfseries\color{navy}}
  {}{0pt}{}
  [\vspace{-4pt}\color{navy}\titlerule]

\titlespacing*{\section}{0pt}{10pt}{6pt}

\newcommand{\entry}[3]{
  \textbf{#1}\hfill{\small #2}\\
  {\emph{#3}}\par
}

\begin{document}

\begin{center}
{\LARGE\bfseries Shu Yiyang}\\[4pt]
{\normalsize 数学専攻・学部学生}\\[2pt]
中国科学技術大学 少年班学院\\[2pt]
中国・合肥
\quad\textbar\quad
\href{mailto:shuyiy@mail.ustc.edu.cn}
{shuyiy@mail.ustc.edu.cn}
\quad\textbar\quad
\href{https://nicholasshu.com}{個人ウェブサイト}
\end{center}

\vspace{-7pt}

\section*{研究関心分野}

微分幾何学、幾何解析学、凸幾何学、および
カケヤ問題とその関連問題。
特に、凸カケヤ集合とその変種、極値的な幾何配置、
ならびに凸性や幾何学的制約のもとでの
面積・直径・周長の相互関係に関心を持つ。

\section*{学歴}

\entry
{中国科学技術大学（USTC）}
{2025年9月～現在}
{少年班学院・数学専攻（純粋数学コース）}

中国科学技術大学少年班・少年創新班の
選抜においてAランクの認定を受け、入学。

\entry
{安徽省安慶第一中学}
{2023年～2025年}
{中国・安徽省安慶市}

2023年および2024年に安徽省の
高校数学コンテストで省二等賞を受賞。
2023年の西安交通大学少年班選抜試験では、
数学科目で全国上位20位以内に入った。

\section*{研究活動}

\entry
{凸カケヤ集合に関する研究}
{2026年9月～現在}
{指導教員：王作勤教授（中国科学技術大学）}

\begin{itemize}
  \item
  あらゆる方向の長さ1の線分を含む凸集合について、
  面積・直径・周長の間の定量的関係を研究している。

  \item
  極値的な幾何配置、幾何不等式、
  および凸体の極限挙動を考察している。

  \item
  支持超平面、支持関数、凸包など、
  凸幾何学の関連手法を研究している。
\end{itemize}

\section*{数学論文・執筆活動}

\textbf{双曲平面のいくつかの幾何学的性質}
\hfill{\small 2026年}\\
\emph{
個人研究論文。
中国科学技術大学数学科学学院の
学生刊行物『蛙鳴』第68号に掲載。
}

\begin{itemize}
  \item
  双曲幾何学のポアンカレ上半平面モデルについて、
  ヒルベルトの公理体系との整合性や
  その変換群などの性質を考察した。

  \item
  双曲三角形の正弦定理・余弦定理、
  およびデザルグの定理、
  パップスの定理などを研究した。
\end{itemize}

\section*{受賞歴・表彰}

\begin{itemize}
\item \textbf{2026年：} 中国科学技術大学 Robogame コンテスト第1位。
\item \textbf{2025年：} CMC 数学コンテスト省二等賞。

  \item
  \textbf{2025年：}
  中国科学技術大学少年創新班入学選抜において
  Aランクの認定を取得。

  \item
  \textbf{2023年・2024年：}
  安徽省高校数学コンテスト
  省二等賞（各年）。

  \item
  \textbf{2023年：}
  西安交通大学少年班選抜試験において、
  数学科目全国上位20位以内。

  \item
  \textbf{2023年：}
  丘成桐中学生数学サマーキャンプを修了。

  \item
  \textbf{2019年・2020年：}
  情報オリンピック関連の省級大会で
  二等賞（各年）。

  \item
  その他、高校在学中に学生自治活動や
  中国共産主義青年団での活動に関する
  複数の表彰を受けた。

\item \textbf{2025年：} 中国科学技術大学「九章杯」三等賞。
\end{itemize}

\section*{数学の学習歴}

\textbf{基礎科目：}
数学解析、線形代数学、常微分方程式論、
古典微分幾何学、凸解析。\\

\textbf{発展的な学習・自主研究：}
滑らかな多様体、微分位相幾何学、
幾何解析学、偏微分方程式論、
極小曲面論、微分形式、
ド・ラームコホモロジー。\\

\emph{
注：上記には大学での履修科目と
自主学習の内容の両方を含む。
個々の科目の成績は記載していない。
}

\section*{学内活動・スキル}

\textbf{学内活動：}
クラス共産主義青年団支部書記
（2025年～現在）；
学生会学術実践部部員
（2026年～現在）。\\

\textbf{技術スキル：}
LaTeX、XeLaTeX、Beamer、SolidWorks。\\

\textbf{言語：}
中国語（母語）；
英語（習熟度未記載）。

\end{document}
